\chapter{$C_2$-Graded Algebras and the Brauer-Wall Group}\label{appendix:gradedalgebras}

In this appendix we collect together some basic results on $C_2$-graded algebras which we make use of in Sections \ref{sect:BrauerWall} and \ref{sect:recoverDyson10}, which we were unable to find in the literature.

We first recall some basic definitions, all of which can be found in \cite[Chap.\ IV]{LAM}. Throughout, $\bK$ is a field with $\charfield(\bK) \neq 2$.

\begin{defn}
    By a \emph{$C_2$-graded algebra} over $\bK$, we mean a finite-dimensional $\bK$-algebra $A$, with $\bK$-subspaces $A_0, A_1 \leq A$, such that $\bK \subset A_0$, $A = A_0 \oplus A_1$, and $A_i \cdot A_j \subset A_{i+j}$, where the subscripts are read in $\bZ / 2 \bZ$.

    A $\bK$-subspace $S \leq A$ of a $C_2$-graded algebra $A$ is called \emph{graded} if $S = S \cap A_0 + S \cap A_1$. 

    A $C_2$-graded algebra $A$ is called \emph{simple} if $A$ has no proper non-zero two-sided graded ideals, and called a \emph{$C_2$-graded division algebra} if every non-zero homogeneous element is invertible.

    The \emph{graded centre} of a $C_2$-graded algebra $A$ is the graded $\bK$-algebra $\widehat{Z}(A)$ such that 
    \begin{align*}
        \widehat{Z}(A)_0 &= \{a \in A_0 \mid ab = ba \text{ for all } b \in A\}, \\
        \widehat{Z}(A)_1 &= \{a \in A_1 \mid ab = ba \text{ for all } b \in A_0, \: ab = - ba \text{ for all } b \in A_1\},
    \end{align*}
and $A$ is called \emph{central} if $\widehat{Z}(A) = \bK$.
\end{defn}

We will use the following basic facts.

\begin{lemma}\label{lem:centreofC2graddiv}
    Suppose that $A$ is a $C_2$-graded simple algebra over $\bK$. Then $\widehat{Z}(A)_1=0$, so
        \[
            \widehat{Z}(A) = A_0 \cap Z(A) \leq Z(A_0).
        \] 
        In particular, if $D$ is a $C_2$-graded division algebra, then $\widehat{Z}(D)$ is a finite field extension of $\bK$.
\end{lemma}

\begin{proof}
    As $\widehat{Z}(A)$ is a graded subalgebra, it suffices to show that $\widehat{Z}(A) \cap A_1 = 0$. Suppose that $z \in \widehat{Z}(A)_1$ and $z \neq 0$. Then $(z)$, the two-sided ideal of $A$ generated by $z$, is homogeneous (as $z$ is homogeneous) and non-trivial, hence $(z) = A$. Furthermore, from the condition that $z$ either commutes or anticommutes with homogeneous elements of $A$, $(z) = Az = \{a z \mid a \in A\}$, so $z$ is invertible. On the other hand, that $z$ is in $\widehat{Z}(A)_1$ means that $z \cdot z = - z \cdot z$, hence $2 z^2 = 0$, so $z^2 = 0$, as $\charfield(\bK) \neq 2$, which contradicts the invertibility of $z$.
\end{proof}

\begin{lemma}\label{lem:invoddeltgradsimple}
    Suppose that $A$ is a $C_2$-graded algebra such that $A_0$ is simple and $A^{\times} \cap A_1 \neq \emptyset$. Then $A$ is simple as a $C_2$-graded algebra.
\end{lemma}

\begin{proof}
    If $I$ is a two-sided graded ideal of $A$, then $I_0$ is a two-sided ideal of $A_0$, hence either $I_0 = 0$ or $I_0 = A_0$. If $u \in A^{\times} \cap A_1$, then left-multiplication by $u$ defines a bijection between $I_0$ and $I_1$, hence either $I = 0$ or $I = A$.
\end{proof}

We now describe simple $C_2$-graded $\bK$-algebras in terms of $C_2$-graded division algebras over $\bK$.

\begin{defn}
    For $p,q \geq 0$ with $p+q \geq 1$, and $A$ a $C_2$-graded $\bK$-algebra, $M_{p,q}(A)$ is $M_{p+q}(A)$ as a $\bK$-algebra, with grading
    \[
    M_{p,q}(A)_0 =
\begin{pmatrix}
M_{p\times p}(A_0) & M_{p\times q}(A_1) \\
M_{q\times p}(A_1) & M_{q\times q}(A_0)
\end{pmatrix},
\qquad
M_{p,q}(A)_1 =
\begin{pmatrix}
M_{p\times p}(A_1) & M_{p\times q}(A_0) \\
M_{q\times p}(A_0) & M_{q\times q}(A_1)
\end{pmatrix}.
\]
\end{defn}

Any central simple $C_2$-graded $\bK$-algebra $A$ is isomorphic to $M_{p,q}(D)$ for some $C_2$-graded division $\bK$-algebra $D$ \cite[Thm.\ 2.10.10]{NO}. To show that this $D$ is unique, we introduce the following.

\begin{defn}
    Suppose that $A$ is a $C_2$-graded $\bK$-algebra. A $C_2$-graded $A$-module $M$ is a $C_2$-graded abelian group $M = M_0 \oplus M_1$, which is also an $A$-module, such that $A_i \cdot M_j \subset M_{i+j}$ (where the subscripts are read in $\bZ / 2 \bZ$). A morphism $f \colon M \rightarrow N$ is $C_2$-graded $A$-modules is an $A$-module homomorphism such that $f(M_i) \subset N_i$ for $i = 0,1$.
    
    We write $\Mod_A^{\wedge}$ for the category of $C_2$-graded $A$-modules.
\end{defn}

The category $\Mod_A^{\wedge}$ has an endofunctor $T \colon \Mod_A^{\wedge} \rightarrow \Mod_A^{\wedge}$ which on objects has $T(M) = M$ as an $A$-module but with the grading $T(M)_0 = M_1$ and $T(M)_1 = M_0$, and is the identity on morphisms.

\begin{defn}
    For $C_2$-graded $\bK$-algebras $A$ and $B$, we say that a functor $F \colon \Mod_A^{\wedge} \rightarrow \Mod_B^{\wedge}$ is \emph{$\bK$-linear} if for any $M,N \in \Mod_A^{\wedge}$, the map
    \[
    F(-) \colon \Hom_A^{\wedge}(M,N) \rightarrow \Hom_B^{\wedge}(F(M), F(N)) 
    \]
    is $\bK$-linear, and \emph{$C_2$-graded} if $F \circ T_A = T_B \circ F$. We say that $\Mod_A^{\wedge}$ and $\Mod_B^{\wedge}$ are equivalent as $\bK$-linear $C_2$-graded categories if there is a $\bK$-linear $C_2$-graded equivalence from $\Mod_A^{\wedge}$ to $\Mod_B^{\wedge}$.
\end{defn}

For a $C_2$-graded $\bK$-algebra $A$, there is an internal $\Hom$ in the category $\Mod_A^{\wedge}$, given by setting 
\[
\underline{\Hom}_A^{\wedge}(M,N) \coloneqq \Hom_A(M,N),
\]
with grading 
\[
\underline{\Hom}_A^{\wedge}(M,N)_0 = \Hom_A^{\wedge}(M,N), \qquad \underline{\Hom}_A^{\wedge}(M,N)_1 = \Hom_A^{\wedge}(M,T(N)).
\]
When $M=N$, this makes $\underline{\End}_A^{\wedge}(M)$ a $C_2$-graded $\bK$-algebra. The following is direct to verify.

\begin{lemma}\label{lem:gradedequivbasicproperties}
    Suppose that $A$ and $B$ are $C_2$-graded $\bK$-algebras. Then 
    \begin{itemize}
        \item For any isomorphism $f \colon A \rightarrow B$ is $C_2$-graded $\bK$-algebras, pushforward
        \[
        f_* \colon \Mod_{B}^{\wedge} \rightarrow \Mod_{A}^{\wedge}
        \] 
        is a $\bK$-linear $C_2$-graded equivalence.
        \item Writing $\! {}_A A$ for $A$ viewed as a left $A$-module and with same grading as $A$, the natural map
        \[
A^{\opp} \rightarrow \underline{\End}_A^{\wedge}({}_A A)
\]
is an isomorphism of $C_2$-graded algebras (where $A^{\opp}$ is the opposite algebra of $A$ (not the graded opposite algebra) and carries the same grading as $A$).
    \end{itemize}
\end{lemma}

We are interested in understanding when two categories $\Mod_{A}^{\wedge}$ and $\Mod_{B}^{\wedge}$ are equivalent. One way to produce functors $\Mod_{A}^{\wedge} \rightarrow \Mod_{B}^{\wedge}$ is through the following construction.

Suppose that $X$ is a $C_2$-graded $(B,A)$-bimodule, meaning that $X$ is a $C_2$-graded left $B$-module and a $C_2$-graded right $A$-module, such that the $C_2$-grading is the same and the actions of $A$ and $B$ commute. For any $M \in \Mod_A^{\wedge}$ we can consider the left $B$-module $X \otimes_A M$, and define
\[
(X \otimes_A M)_i = \langle x \otimes m \mid x \in X, m \in M \text{ homogeneous with } \deg(x) + \deg(m) = i\rangle_{\bZ},
\]
which gives $X \otimes_A M$ the structure of a $C_2$-graded $B$-module \cite[\S 2.4]{NO}. The resulting functor
\[
X \otimes_A (-) \colon \Mod_{A}^{\wedge} \rightarrow \Mod_{B}^{\wedge}
\]
is $C_2$-graded, and $\bK$-linear whenever the $\bK$-linear structures on $X$ from $A$ and $B$ agree.

For the pair $A, M_{p,q}(A)$, with $n \coloneqq p+q \geq 1$, we consider the $C_2$-graded $(A, M_{p,q}(A))$-bimodule
\[
A^{1 \times n}_{p,q}, \: \text{ with } \: (A^{1 \times n}_{p,q})_0 = A_0^{1 \times p} \oplus A_1^{1 \times q}, \: \text{ and } \: (A^{1 \times n}_{p,q})_1 = A_1^{1 \times p} \oplus A_0^{1 \times q}
\]
of row vectors, and the $C_2$-graded $(M_{p,q}(A), A)$-bimodule of column vectors
\[
A^{n \times 1}_{p,q}, \: \text{ with } \: (A^{n \times 1}_{p,q})_0 = A_0^{p \times 1} \oplus A_1^{q \times 1}, \: \text{ and } \: (A^{n \times 1}_{p,q})_1 = A_1^{p \times 1} \oplus A_0^{q \times 1}.
\]

\begin{lemma}\label{lem:gradedmorita}
    Let $A$ be a $C_2$-graded $\bK$-algebra and $p,q \geq 0$ with $n \coloneqq p+q \geq 1$. Then the functors 
        \[
        F \colon \Mod_A^{\wedge} \rightarrow \Mod_{M_{p,q}(A)}^{\wedge}, \qquad F(M) \coloneqq A^{1 \times n}_{p,q} \otimes_{M_{p,q}(A)} M,
        \]
        and
        \[
        G \colon \Mod_{M_{p,q}(A)}^{\wedge} \rightarrow \Mod_A^{\wedge}, \qquad G(N) \coloneqq A^{n \times 1}_{p,q} \otimes_{A} N,
        \]
        are mutually inverse equivalences of $\bK$-linear $C_2$-graded categories.
\end{lemma}

\begin{proof}
    From the above discussion, both functors are $C_2$-graded, and $\bK$-linear because the two possible $\bK$-linear structures on each of the bimodules $A^{n \times 1}_{p,q}$ and $A^{1 \times n}_{p,q}$ coincide. To see that they are mutually equivalent, it suffices to show that there are isomorphisms
    \[
        \phi \colon A^{n \times 1}_{p,q} \otimes_A A^{1 \times n}_{p,q} \xrightarrow{\sim} M_{p,q}(A), \qquad \psi \colon A^{1 \times n}_{p,q} \otimes_A A^{n \times 1}_{p,q} \xrightarrow{\sim} A
    \]
    of $C_2$-graded $(M_{p,q}(A), M_{p,q}(A))$-bimodules and $(A,A)$-bimodules respectively. Let $e_1, ... , e_n$ denote the standard $A$-basis of $A^{n \times 1}_{p,q}$ and $f_1, \cdots, f_n$ the standard $A$-basis of $A^{1 \times n}_{p,q}$. We define $\phi$ and $\psi$ by
    \[
        \phi\left(\sum_{i,j} a_{ij} e_i \otimes f_j\right) = \sum_{i,j} a_{ij} E_{ij}, \qquad  \psi\left(\sum_{i,j} a_{ij} f_i \otimes e_j\right) = \sum_{i} a_{ii}.
    \]
    It is direct to check that these are bimodule isomorphisms, which respect the $C_2$-grading.
    \end{proof}

\begin{lemma}\label{lem:divalgcat}
    Suppose that $D$ is a $C_2$-graded division algebra over $\bK$. Then
    \begin{itemize}
        \item If $D_1 \neq 0$, every object of $\Mod_{D}^{\wedge}$ is isomorphic to a direct sum of copies of ${}_D D$.
        \item If $D_1 = 0$, then object of $\Mod_{D}^{\wedge}$ is isomorphic to a direct sum of copies of ${}_D D$ and $T({}_D D)$.
    \end{itemize}
\end{lemma}

\begin{proof}
    Suppose that $D_1 \neq 0$ and $M$ is an object of $\Mod_D^{\wedge}$. $D_0$ is a $\bK$-division algebra, so there is a basis $\{e_i\}_{i \in I}$ of $M_0$ as a $D_0$-module. It is straightforward to show that
    \[
    \phi \colon \bigoplus_{i \in I} {}_D D \rightarrow M, \qquad \phi\left(\sum_i x_i\right) \coloneqq x_i e_i
    \]
    is an isomorphism in $\Mod_D^{\wedge}$, using that multiplication by any $x \in D_1 \setminus 0$ defines a $\bK$-linear isomorphism from $M_0$ to $M_1$. 

    Now suppose that $D_1 = 0$, and $M$ an object of $\Mod_D^{\wedge}$. We may similarly take a basis $\{e_i\}_{i \in I}$ of $M_0$ and $\{f_j\}_{j \in J}$ of $M_1$ as $D_0$-modules, and define
    \[
    \psi \colon \bigoplus_{i \in I} {}_D D \oplus \bigoplus_{j \in J} T({}_D D) \rightarrow M, \qquad \phi\left(\sum_i x_i, \sum_j y_j \right) \coloneqq \sum_i x_i e_i + \sum_j y_j f_j,
    \]
    which is also easily seen to be an isomorphism in $\Mod_D^{\wedge}$.
\end{proof}

\begin{lemma}\label{lem:recovergradeddivisionalgebra}
    Suppose that $D$ and $D'$ are $C_2$-graded division algebras over $\bK$. If $\Mod_{D}^{\wedge}$ and $\Mod_{D'}^{\wedge}$ are equivalent as $\bK$-linear $C_2$-graded categories, then $D_A \cong D_B$ as $C_2$-graded $\bK$-algebras.
\end{lemma}

\begin{proof}
    Let $F \colon \Mod_{D'}^{\wedge} \rightarrow \Mod_{D}^{\wedge}$ be such an equivalence, and set $X \coloneqq {}_{D'} D'$. The object $Y \coloneqq F(X) \in \Mod_{D}^{\wedge}$ is simple, and therefore isomorphic to ${}_D D$ or $T({}_D D)$ by Lemma \ref{lem:divalgcat}. The functor $T \colon \Mod_{D}^{\wedge} \rightarrow \Mod_{D}^{\wedge}$ is a $\bK$-linear $C_2$-graded equivalence, and therefore potentially replacing $F$ with $T \circ F$ we may assume that $Y \cong {}_D D$. Using that $F(T'(X)) = T(F(X)) = T(Y)$, the functor $F$ induces identifications
    % https://q.uiver.app/#q=WzAsNixbMCwwLCJcXHVuZGVybGluZXtcXEVuZH1fe0QnfV57XFx3ZWRnZX0oWCkiXSxbMSwwLCJcXEhvbV97RCd9XntcXHdlZGdlfShYLFgpIl0sWzIsMCwiXFxIb21fe0QnfV57XFx3ZWRnZX0oWCxUJyhYKSkiXSxbMCwxLCJcXHVuZGVybGluZXtcXEVuZH1fRF57XFx3ZWRnZX0oWSkiXSxbMSwxLCJcXEhvbV97RH1ee1xcd2VkZ2V9KFksWSkiXSxbMiwxLCJcXEhvbV97RH1ee1xcd2VkZ2V9KFksVChZKSkiXSxbMCwxLCI9IiwxLHsic3R5bGUiOnsiYm9keSI6eyJuYW1lIjoibm9uZSJ9LCJoZWFkIjp7Im5hbWUiOiJub25lIn19fV0sWzMsNCwiPSIsMSx7InN0eWxlIjp7ImJvZHkiOnsibmFtZSI6Im5vbmUifSwiaGVhZCI6eyJuYW1lIjoibm9uZSJ9fX1dLFs0LDUsIlxcYmlnb3BsdXMiLDEseyJzdHlsZSI6eyJib2R5Ijp7Im5hbWUiOiJub25lIn0sImhlYWQiOnsibmFtZSI6Im5vbmUifX19XSxbMSwyLCJcXGJpZ29wbHVzIiwxLHsic3R5bGUiOnsiYm9keSI6eyJuYW1lIjoibm9uZSJ9LCJoZWFkIjp7Im5hbWUiOiJub25lIn19fV0sWzEsNCwiRigtKSIsMl0sWzIsNSwiRigtKSIsMl1d
\[\begin{tikzcd}
	{\underline{\End}_{D'}^{\wedge}(X)} & {\Hom_{D'}^{\wedge}(X,X)} & {\Hom_{D'}^{\wedge}(X,T'(X))} \\
	{\underline{\End}_D^{\wedge}(Y)} & {\Hom_{D}^{\wedge}(Y,Y)} & {\Hom_{D}^{\wedge}(Y,T(Y))}
	\arrow["{=}"{description}, draw=none, from=1-1, to=1-2]
	\arrow["\bigoplus"{description}, draw=none, from=1-2, to=1-3]
	\arrow["{F(-)}"', from=1-2, to=2-2]
	\arrow["{F(-)}"', from=1-3, to=2-3]
	\arrow["{=}"{description}, draw=none, from=2-1, to=2-2]
	\arrow["\bigoplus"{description}, draw=none, from=2-2, to=2-3]
\end{tikzcd}\]
The equivalence $F$ is between additive categories, hence $F$ is additive, and hence the induced identification $\Phi \colon \underline{\End}_{D'}^{\wedge}(X) \rightarrow \underline{\End}_D^{\wedge}(Y)$ commutes with addition. For multiplication,
\begin{align*}
    \Phi(f_0, f_1) \cdot \Phi(g_0, g_1) &= (F(f_0), F(f_1)) \cdot (F(g_0), F(g_1)), \\
    &= (F(f_0) \circ F(g_0) + F(f_1) \circ F(g_1), F(f_0) \circ F(g_1) + F(f_1) \circ F(g_0)), \\
    &= (F(f_0 \circ g_0) + F(f_1 \circ g_1), F(f_0 \circ g_1) + F(f_1 \circ g_0)), \\
    &= \Phi((f_0, f_1) \cdot (g_0,g_1)),
\end{align*}
again using that $F$ is additive. It is also easy to see that $\Phi$ preserves the identity, and therefore the $\bK$-algebra structure because $F$ is $\bK$-linear. Using Lemma \ref{lem:gradedequivbasicproperties}, we therefore have isomorphisms
\[
(D')^{\opp} \cong \underline{\End}_{D'}^{\wedge}({}_{D'}D') = \underline{\End}_{D'}^{\wedge}(X) \xrightarrow{\Phi} \underline{\End}_D^{\wedge}(Y) \cong \underline{\End}_{D}^{\wedge}({}_D D) \cong D^{\opp}
\]
of $C_2$-graded $\bK$-algebras, and thus $D \cong D'$ as $C_2$-graded $\bK$-algebras.
\end{proof}

\begin{cor}\label{cor:stdformforgradedalg}
If $A$ is a simple $C_2$-graded algebra over $\bK$, then 
    \[
    A \cong M_{p,q}(D_A)
    \]
for a unique $C_2$-graded division algebra $D_A$ over $\bK$.
\end{cor}

\begin{proof}
    The existence is \cite[Thm.\ 2.10.10]{NO}. For uniqueness, if additionally $A \cong M_{p',q'}(D')$ for some other $C_2$-graded division algebra $D'$ over $\bK$, then $\Mod_{D}^{\wedge}$ and $\Mod_{D'}^{\wedge}$ are equivalent as $\bK$-linear $C_2$-graded categories by Lemma \ref{lem:gradedequivbasicproperties} and Lemma \ref{lem:gradedmorita}, and hence $D \cong D'$ as $C_2$-graded $\bK$-algebras by Lemma \ref{lem:recovergradeddivisionalgebra}.
\end{proof}

\begin{cor}\label{cor:morita}
Let $A$ and $B$ be simple $C_2$-graded $\bK$-algebras. Then $\Mod_A^{\wedge}$ and $\Mod_B^{\wedge}$ are equivalent as $\bK$-linear $C_2$-graded categories if and only if $D_A \cong D_B$ as $C_2$-graded $\bK$-algebras.
\end{cor}

\begin{proof}
    If $D_A \cong D_B$ as $C_2$-graded $\bK$-algebras, then 
    \[
    \Mod_A^{\wedge} \cong \Mod_{D_A}^{\wedge} \cong \Mod_{D_B}^{\wedge} \cong \Mod_B^{\wedge}
    \]
    $\bK$-linear $C_2$-graded categories by Lemma \ref{lem:gradedequivbasicproperties} and Lemma \ref{lem:gradedmorita}. Conversely, if $\Mod_A^{\wedge}$ and $\Mod_B^{\wedge}$ are equivalent as $\bK$-linear $C_2$-graded categories, then so are $\Mod_{D_A}^{\wedge}$ and $\Mod_{D_B}^{\wedge}$ by Lemma \ref{lem:gradedmorita}, and thus $D_A \cong D_B$ as $C_2$-graded $\bK$-algebras by Lemma \ref{lem:recovergradeddivisionalgebra}.
\end{proof}

\begin{remark}
    We note that in all the above one could replace $\Mod_A^{\wedge}$ with the full subcategory of \emph{finite-dimensional} $C_2$-graded $A$-modules, and all statements and proofs remain valid.
\end{remark}

We can use the above results to prove a $C_2$-graded double centraliser theorem. We recall \cite[\S IV.2]{LAM} that for a $C_2$-graded $\bK$-algebra $A$ and graded subspace $S \subset A$, the graded centraliser is the graded subalgebra $\widehat{C}_A(S)$ with
\begin{align*}
        \widehat{C}_A(S)_0 &= \{a \in A_0 \mid ab = ba \text{ for all } b \in S\}, \\
        \widehat{C}_A(S)_1 &= \{a \in A_1 \mid ab = ba \text{ for all } b \in S_0, \: ab = - ba \text{ for all } b \in S_1\}.
    \end{align*}
For example, $\widehat{Z}(A) = \widehat{C}_A(A)$.
\begin{prop}\label{prop:graddoublecent}
    Suppose that $A$ is a central simple $C_2$-graded algebra over $\bK$, and let $B$ be a simple $C_2$-graded subalgebra of $A$. Then
    \begin{itemize}
        \item $C \coloneqq \widehat{C}_A(B)$ is simple as a $C_2$-graded algebra over $\bK$,
        \item $B = \widehat{C}_A(C)$,
        \item $\dim_{\bK} A = \dim_{\bK} B \cdot \dim_{\bK} C$.
    \end{itemize}

    Furthermore, $\widehat{Z}(B) = \bK$ and and only if $\widehat{Z}(C) = \bK$, in which case the natural map
    \[
    B \, \hat{\otimes} \, C \rightarrow A
    \]
    is an isomorphism of $C_2$-graded algebras over $\bK$.
\end{prop}

\begin{proof}
    Write $A^*$ for the graded opposite algebra of $A$ \cite[\S IV.4]{LAM}, and $A^{\opp}$ for the usual opposite algebra. From the proof of \cite[Prop.\ IV.4.3]{LAM} there is an isomorphism
    \[
    \Theta \colon A \hat{\otimes} A^* \xrightarrow{\sim} \End_{\bK}(A)
    \]
    of $C_2$-graded $\bK$-algebras which on homogeneous elements is given by
    \[
    \Theta(b \otimes a^*)(x) = (-1)^{|a||x|} bxa.
    \]
    We can restrict this left $A \hat{\otimes} A^*$-module structure on $A$ to $A^*$, and it is straightforward to show that
    \[
        f_{(-)} \colon (A^*)^{\opp} \rightarrow \underline{\End}_{A^*}^{\wedge}(A), \qquad f_a(x) \coloneqq (-1)^{|a||x|}a x
    \]
    is an isomorphism of $C_2$-graded $\bK$-algebras. Setting $E \coloneqq B \hat{\otimes} A^*$, a direct computation shows that this restricts to a $C_2$-graded $\bK$-algebra isomorphism
    \[
        f_{(-)} \colon (\widehat{C}_A(B)^*)^{\opp} \xrightarrow{\sim} \underline{\End}_{E}^{\wedge}(A).
    \]

    By \cite[Thm.\ IV.2.3]{LAM}, $E$ is a simple $C_2$-graded algebra over $\bK$, and therefore $E \cong M_{p,q}(D_E)$ by Corollary \ref{cor:stdformforgradedalg}. From Lemma \ref{lem:gradedmorita}, there is an equivalence $F \colon \Mod_{D_E}^{\wedge} \xrightarrow{\sim} \Mod_E^{\wedge}$, and therefore by Lemma \ref{lem:divalgcat} any object of $\Mod_E^{\wedge}$ is isomorphic to a direct sum of copies of $X \coloneqq F({}_{D_E} D_E)$ and $T(X)$ (which themselves are isomorphic if and only if $(D_E)_1 \neq 0$). In particular, $A \cong X^{\oplus m_1} \oplus T(X)^{\oplus m_2}$ for some $m_1, m_2 \geq 0$. We have that $D_E^{\opp} \xrightarrow{\sim} \underline{\End}_E^{\wedge}(X)$ by as $C_2$-graded $\bK$-algebras by Lemma \ref{lem:gradedequivbasicproperties}, and
    \[
    T(-) \colon \underline{\End}_E^{\wedge}(X) \xrightarrow{\sim} \underline{\End}_E^{\wedge}(T(X))
    \]
    is an isomorphism of $C_2$-graded $\bK$-algebras (following the same reasoning as in the proof of Lemma \ref{lem:recovergradeddivisionalgebra}). Furthermore, the natural map
    \[
    M_{m_1, m_2}(\underline{\End}_E^{\wedge}(X)) \rightarrow \underline{\End}_E^{\wedge}(X^{\oplus m_1} \oplus T(X)^{\oplus m_2})
    \]
    is an isomorphism of $C_2$-graded $\bK$-algebras, and therefore
    \[
        M_{m_1, m_2}(D_E^{\opp}) \cong \underline{\End}_E^{\wedge}(X^{\oplus m_1} \oplus T(X)^{\oplus m_2}) \cong \underline{\End}_E^{\wedge}(A) \cong (\widehat{C}_A(B)^*)^{\opp}
    \]
    as $C_2$-graded $\bK$-algebras. In particular, $C = \widehat{C}_A(B)$ is simple as a $C_2$-graded $\bK$-algebra.

    We have that $B \subset \widehat{C}_A(C)$, and therefore to see that $B \subset \widehat{C}_A(C)$ it suffices to show the third claim that $\dim_{\bK} A = \dim_{\bK} B \cdot \dim_{\bK} C$. Let us denote these quantities as $a,b$ and $c$, and similarly for $d, e, x$. 

    By definition, $e = ab$, and setting $m \coloneqq m_1 + m_2$, $a = m x$ and $c = m^2 d$ from the isomorphism $(\widehat{C}_A(B)^*)^{\opp} \cong M_{m_1, m_2}(D_E^{\opp})$. Setting $n \coloneqq p+q$, then from the isomorphism $E \cong M_{p,q}(D_E)$, $e = n^2 d$ and $x = n d$ from the definition of the equivalence of Lemma \ref{lem:gradedmorita}. Together,
    \[
    ab = e = n^2 d = \left(\frac{x}{d} \right)^2 d = \frac{x^2}{d} = \left(\frac{a}{m} \right)^2 \frac{1}{d}= \frac{a^2}{m^2 d} = \frac{a^2}{c},
    \]
    and so $a = bc$ as required.
    
    Finally, note that if $\widehat{Z}(B) = \bK$, then $B \hat{\otimes} C$ is simple by \cite[Thm.\ IV.2.3]{LAM}, and therefore the natural map $B \, \hat{\otimes} \, C \rightarrow A$ has trivial kernel, and thus is an isomorphism by comparing the dimensions. As this induces an isomorphism $\widehat{Z}(B) \hat{\otimes} \widehat{Z}(C) \xrightarrow{\sim} \widehat{Z}(A)$ \cite[Thm.\ IV.2.3]{LAM}, it follows that $\widehat{Z}(C) = \bK$. The same argument with the roles of $B$ and $C$ reverses shows that $\widehat{Z}(C) = \bK$ implies that $\widehat{Z}(B) = \bK$.
\end{proof}

The above results have consequences for the Brauer-Wall group \cite[\S IV.4]{LAM}.

\begin{defn}\label{def:BW}
    The Brauer-Wall group of $\bK$, $\BrW(\bK)$, is the set of central simple $C_2$-graded algebras over $\bK$ up to equivalence, where $[A] = [B]$ in $\BrW(\bK)$ if and only if $M_{p,q}(A) \cong M_{r,s}(B)$ for some $p,q,r,s \geq 0$ with $p+q, r+s \geq 1$. 
    
    The multiplication is defined by the graded tensor product
    \[
    [A] \cdot [B] \coloneqq [A \hat{\otimes} B].
    \]
\end{defn}

This is related to the classical Brauer group as follows. First, viewing any simple $\bF$-algebra $A$ as a simple $C_2$-graded $\bF$-algebra with the trivial grading defines a homomorphism
\[
\Br(\bF) \rightarrow \BrW(\bF).
\]
This forms the first term of the exact sequence \cite[Thm.\ IV.4.4]{LAM}
\begin{equation}\label{eqn:exactseqbrauer}
0 \rightarrow \Br(\bF) \rightarrow \BrW(\bF) \xrightarrow{j} Q(\bF) \rightarrow 0,
\end{equation}
where 
\[
Q(\bF) \coloneqq \bZ / 2 \bZ \: \dot{\times} \: \bF^{\times} / \bF^{\times 2},
\]
the direct product as a set, but with twisted multiplication 
\[
(e,r) * (e', r') = (e+e', (-1)^{ee'} rr').
\]
The group homomorphism $j \colon \BrW(\bF) \rightarrow Q(\bF)$ is defined in \cite[Thm.\ IV.4.3]{LAM}.

We write $\GradDiv_{\bK}^{C_2}$ for the set of \emph{central} $C_2$-graded division $\bK$-algebras up to isomorphism. The following shows that we can think of $\BrW(\bK)$ is defining a group structure on this set.

\begin{cor}\label{cor:divalgrep}
    The map
    \[
    \GradDiv_{\bK}^{C_2} \rightarrow \BrW(\bK), \qquad D \mapsto [D]
    \]
    is an isomorphism: every class in $\BrW(\bK)$ has a unique representative in $\GradDiv_{\bK}^{C_2}$.
\end{cor}

\begin{proof}
    There is a canonical $\bK$-algebra identification $\widehat{Z}(D) \rightarrow \widehat{Z}(M_{p,q}(D))$ defined by
    \[
    z_0 + z_1 \mapsto \begin{pmatrix}
        (z_0 + z_1) I_p & 0_{p,q} \\ 0_{q,p} & (z_0 - z_1) I_q
    \end{pmatrix},
    \]
    and therefore the claim follows directly from Corollary \ref{cor:stdformforgradedalg}.
\end{proof}

\begin{cor}
    If $A$ and $B$ are central simple $C_2$-graded $\bK$-algebras, then $[A] = [B]$ in $\BrW(\bK)$ if and only if $\Mod_A^{\wedge}$ and $\Mod_B^{\wedge}$ are equivalent as $\bK$-linear $C_2$-graded categories.
\end{cor}

\begin{proof}
    This follows directly from Corollary \ref{cor:morita} and Corollary \ref{cor:divalgrep}.
\end{proof}



